% Generated draft: requires mathematical and source review.
\documentclass[11pt]{article}
\usepackage[margin=25mm]{geometry}
\usepackage{fontspec}
\setmainfont{Latin Modern Roman}
\usepackage{amsmath,amssymb,mathtools,unicode-math}
\setmathfont{Latin Modern Math}
\usepackage{enumitem,needspace,xurl,hyperref}
\setlength{\parindent}{0pt}
\setlength{\parskip}{5pt}
\setlength{\emergencystretch}{3em}
\newcommand{\N}{\mathbb{N}}
\newcommand{\Z}{\mathbb{Z}}
\newcommand{\Q}{\mathbb{Q}}
\newcommand{\R}{\mathbb{R}}
\newcommand{\C}{\mathbb{C}}
\newcommand{\F}{\mathbb{F}}
\newcommand{\A}{\mathbb{A}}
\newcommand{\PP}{\mathbb{P}}
\newcommand{\E}{\mathbb{E}}
\newcommand{\eps}{\varepsilon}
\newcommand{\dd}{\,\mathrm{d}}
\begin{document}
\section*{UnsolvedMath ID: 30000663}
\textbf{Source problem code:} catalogue entry
\subsection{Definitions and mathematical statement}
Let $X=\operatorname{Spec}A$ be a normal affine complex variety, let an algebraic group $G$ act effectively on $X$, and let $\varphi:X\to X$ be an unramified endomorphism satisfying
\[
\varphi(gx)=g\varphi(x)\qquad(g\in G,\ x\in X).
\]
Effectivity means that only the identity acts trivially on all of $X$. Put $B=A^G$. Equivariance implies that $\varphi^*:A\to A$ restricts to an endomorphism $\varphi^*|_B:B\to B$, giving
\[
\psi:\operatorname{Spec}B\longrightarrow\operatorname{Spec}B.
\]
With $X//G=\operatorname{Spec}A^G$ as the affine invariant-theoretic quotient, the question is whether $\psi$ must be unramified.

For a morphism of schemes, unramified means locally of finite type and with zero sheaf of relative Kähler differentials. Thus the ring formulation asks whether the induced endomorphism of $B$ has these properties, given them for $\varphi$. In the finite-type-variety setting, the finite-type condition is automatic for an endomorphism, and the differential condition is $\Omega_{B/B,\,\varphi^*|_B}=0$, with the two copies of $B$ related by the displayed homomorphism.

The source permits an arbitrary algebraic group and does not assume reductivity or finite generation of $A^G$. Consequently its quotient need not automatically be an affine variety of finite type. The affine-scheme interpretation above preserves that group scope; a version restricted to finite-type quotients must state the additional hypothesis. No smoothness of $X$, flatness of $\varphi$, or reductivity of $G$ is silently added.
\subsection{Short English statement}
Does an equivariant unramified self-map of a normal affine complex variety induce an unramified self-map on its invariant quotient? The invariant-ring construction is explicit; arbitrary algebraic groups require care because the quotient ring need not be finitely generated.
\subsection{Sources}
\begin{itemize}
\item Supplied report: \emph{Affine Algebraic Geometry} (2007), Oberwolfach Reports, DOI 10.4171/owr/2007/01, pp. 5--82; proposed year 2007. \url{https://doi.org/10.4171/owr/2007/01}.
\item Local UnsolvedMath registry, permanent ID 30000663, catalogue entry; created 25 August 2026, updated 16 September 2026. Its corrected\_verified note describes balancing notation. The supplied 21 August 2026 literature triage labels the extracted claim open and reports no located authoritative complete resolution. No additional external reference is supplied, and no independent current-status assessment is asserted here.
\end{itemize}
\end{document}
